\section{如何学好刑诉法}
\begin{enumerate}
  \item 做题。
        \begin{enumerate}
          \item 抓关键字。\al{在审判阶段，\zd{值班律师}可以出庭进行\zd{辩论}}。
          \item 区分不同主体，不同阶段。
                \begin{enumerate}
                  \item \al{县公安局以甲涉嫌过失致人死亡\xhx{立案侦查}，后随即\xhx{撤销案件}，乙的近亲属的诉讼权利。\scx{就撤销向县公安局复议}。只能公转自了}。
                  \item 注意执行阶段，已无嫌疑人、被告人身份。
                  \item 不要忘了最高法、最高检。
                \end{enumerate}
          \item 以偏概全。\al{外籍甲只能委托具有中国律师资格并取得执业证书的律师作为辩护人}。做梗概。
                \begin{enumerate}
                  \item \al{只能委托律师作为辩护人}。
                  \item \al{仅仅、都要、一律、必须}。
                \end{enumerate}
          \item 选项看起来冠冕堂皇很有理由，应当就是应当，没有例外，一般应当才有例外。
                \begin{enumerate}
                  \item \al{速裁程序，牟县法院经过审理认为案件复杂，可以在庭后听取控辩双方意见后定期宣判}。
                \end{enumerate}
          \item 正话反说、张冠李戴、答非所问/无关选项、声东击西、无中生有、故弄玄虚。
        \end{enumerate}
  \item 7 月份进入了法考备战的黄金时间，需要\zd{对重点的考点进行专项的背诵记忆}。
  \item 客观题对考点的\zd{归纳总结背诵}，对选项的\zd{判断}能力。
  \item 主观题考的是对案情的\zd{阅读}能力，对考点的\zd{识别}能力，法条\zd{搜索}能力，结论的\zd{书写}能力，\zd{论证说理}的能力。
  \item 案例化，法条化，实务操作能力，关注身边的真实的法律新闻和案例。
  \item 关于未成年人犯罪记录封存的实施办法；取保候审的规定；关于办理信息网络案件适用刑事诉讼程序若干问题的意见。
  \item \zd{新法必考，新增必考，重者恒重，热点必考（热点案例/改革热点）}
  \item 上诉不加刑原则，管辖制度，证据制度，辩护制度，认罪认罚从宽制度，附条件不起诉制度。
  \item 重要的法律文件。
        \begin{enumerate}
          \item 刑事诉讼法。
          \item 最高院的刑诉解释、最高检的高检规则、公安部的公安部规定。
          \item 专项问题找专规。排非规定、认罪认罚指导意见。
          \item 2023 年。
                \begin{enumerate}
                  \item 未成年人犯罪记录封存的实施办法。
                  \item 取保候审的若干问题的规定。
                  \item 关于办理信息网络犯罪案件适用刑事诉讼若干问题的规定。
                \end{enumerate}
        \end{enumerate}
  \item 改革热点。
        \begin{enumerate}
          \item 以审判为中心的诉讼制度改革。
          \item 控辩式的审判模式改革。
          \item 认罪认罚从宽制度改革。
          \item 缺席审判改革。
        \end{enumerate}
  \item \zd{全局框架}。三段：
        \begin{enumerate}
          \item 10 项原理和 10 项制度。基本原则、专门机关和诉讼参加人、管辖、回避、\zd{辩护}和代理、证据、\zd{强制措施}、\zd{附带民事}、期间送达。
          \item 立案、侦查、起诉、审判（一审、二审、复核、再审）、执行。
          \item 5 个特别程序加 1 个涉外程序。未成年人、公诉和解、缺席审判、没收、强制医疗、涉外刑诉。
        \end{enumerate}
\end{enumerate}

\section{概说和基本理念}
\begin{enumerate}
  \item 一个概念、两对关系、三组理念和六个范畴。
  \item 刑事诉讼是指在\zd{公检法}\zd{专门机关}的主持下在\zd{当事人和其他诉讼参与人}的参加下，严格按照\zd{法定程序}，解决被追诉者的刑事责任问题的活动。
        \begin{enumerate}
          \item \zd{3 大 4 小来主持}。
          \item \zd{诉讼参与人参加} = 当事人 + 其他诉讼参与人。
          \item 程序。立案 \yjt 侦查 \yjt 起诉 \yjt 审判 \yjt 执行。
        \end{enumerate}
  \item 刑事诉讼法。国家制定和认可的调整刑诉活动的法律规范的总称。有狭义（刑法典）和广义（调整刑诉活动的一切法律规范）之分。
  \item 渊源。在哪些文件中能找到刑诉法的依据。判断标准：\zd{以下 7 个 + 内容相关}。
        \begin{enumerate}
          \item 宪法、刑诉法典。
          \item 法律和立法解释、法律解释和规定、法规和部门规章。
          \item 国际条约和公约。
        \end{enumerate}
  \item 刑诉法和刑法的关系。
        \begin{enumerate}
          \item \zd{工具价值}。为刑法\zd{提供}组织上、基本构架、程序证据规范、标准、准确性、效率性上的\zd{保障}。
          \item \zd{独立价值}。\zd{本身}体现民主法制；\zd{弥补}刑法不足/空白；\zd{影响}刑法结果实现。
        \end{enumerate}
  \item 与法治国家的关系。宪法是\zd{静态}的刑事诉讼法，刑事诉讼法是\zd{动态}的宪法。宪法不会自我实现它的内容，刑事诉讼法进行维护捍卫。
  \item \zd{基本理念}。刑事诉讼的灵魂。整个刑诉流程以及各项制度设计的指导思想。
        \begin{enumerate}
          \item 惩罚犯罪。\zd{打击犯罪}。
          \item 保障人权。
                \begin{enumerate}
                  \item 无辜的人不受追究；有罪的人受到公正处罚；诉讼权利得到充分保障和行使。
                  \item 制度涉及。辩护权、法院援助、非法证据排除。
                \end{enumerate}
          \item \zd{保障人权和惩罚犯罪既对立又统一，二者要并重}。
          \item 实体公正。\zd{结果公正}。事实清楚、证据充分；定罪准确、量刑适当；涉案财物处置公正、有错必纠。
          \item 程序公正。\zd{过程公正}。程序要遵守；权利要保障；措施要适当；刑讯要严禁；程序要透明；审判要中立、控审要分离；期限要及时。
          \item \zd{坚持实体公正和程序公正并重原则}。
          \item 司法公正。公正是\zd{首要目标价值}。
          \item 诉讼效率。简易程序、速裁程序、缺席审判、认罪认罚。
          \item \zd{公正优先，兼顾效率}。
        \end{enumerate}
\end{enumerate}
